\section{A finite description of a minimum lattice}\label{sec:minimum}

The computable ceiling $M_*$ from \Cref{q:size-ceiling} can be combined
with one ordinary Busy Beaver constant to specify a nonrepresentable
lattice of minimum cardinality.  The constant is defined independently
of lattices and representations.  Once it is given, the carrier and
every order-table bit are specified by bounded finite arithmetic.
All arithmetic and all assertions of termination in this section refer
to the ordinary standard natural numbers.

\subsection{The register language and its constant}

\begin{definition}[Register programs and size]\label{q:register-language}
Registers $v_i$, $i\ge0$, hold natural numbers and initially all contain
zero.  Terms use these registers, the constants $0,1$, addition,
truncated subtraction $u\mathbin{\dot-}v$, multiplication, exponentiation
$u^v$, and integer floor quotient.  We set $0^0=1$ and give quotient by
zero the value zero.  Bounded formulas use the binary predicates $=,<$,
the connectives $\neg,\land,\lor$, and quantifiers
$\exists v_i<t$, $\forall v_i<t$, with $v_i$ absent from the bound term
$t$.  A variable bound within a formula is local to that formula and
does not write to a register.

Programs are formed by assignments of terms to registers, binary
sequencing, if-then-else with a bounded condition and two program
branches, and while with a bounded condition and a program body.  They
have the usual deterministic semantics.  A terminating program outputs
its final value of $v_0$.

The size of a program is its syntax-tree cost: each operation, constant,
predicate, connective, quantifier, assignment, sequencing,
if-then-else, or while node costs one.  Each occurrence of $v_i$ costs
$1+i$, including assignment targets and quantifier binders.  Parentheses
cost nothing, and no extra cost is charged for the use of a quantifier's
term as its bound.
\end{definition}

\begin{definition}[The Busy Beaver cutoff]\label{q:BB}
With $M_*=x_5$ as in \eqref{q:x-sequence}, set
\begin{equation}\label{q:BB-bound}
 \begin{split}
 k&=2^{(M_*+10)^4},\\
 B&=1+\max\Bigl(\{0\}\cup
       \{\operatorname{output}(p):p\text{ terminates and has size }
                                      \le k\}\Bigr).
 \end{split}
\end{equation}
The programs in this definition use exactly
\Cref{q:register-language}; in particular $B$ is not a symbol in their
language.
\end{definition}

This is a maximum-output variant of Rad\'o's finite-program extremal
construction \citep{Rado1962}, with the register syntax and cost fixed
here.  It is not the classical two-symbol state-count function.
There are finitely many programs of bounded size, since the cost of a
register occurrence bounds its index.  Thus $B$ is a definite finite
integer.  It is an ordinary, unevaluated Busy Beaver constant, defined
by termination and output of programs without an oracle.

\subsection{Bounded arithmetic for lattice tables and witnesses}

\begin{definition}[Digits and order tables]\label{q:table-predicate}
For $b\ge1$, put
\begin{equation}\label{q:digit}
 \operatorname{dig}_b(z,i)=
   \left\lfloor\frac{z}{b^i}\right\rfloor
       \mathbin{\dot-}
       b\left\lfloor\frac{z}{b^{i+1}}\right\rfloor.
\end{equation}
Base $1$ is used only where the intended digit is zero, in which case
this formula indeed gives zero.  For $n\ge1$ and $0\le t<2^{n^2}$,
let $T_{n,t}$ be the relation on $\{0,\ldots,n-1\}$ defined by
\begin{equation}\label{q:order-bit}
          a\preceq b\quad\Longleftrightarrow\quad
                         \operatorname{dig}_2(t,an+b)=1.
\end{equation}
Write $E(a,b)$ for the right side, with $n,t$ understood.  The bounded
predicate $P(n,t)$ says that $T_{n,t}$ is a lattice: it is the conjunction
of
\begin{align*}
 &\forall a<n\ E(a,a),\\
 &\forall a,b<n\ \bigl((E(a,b)\land E(b,a))\Longrightarrow a=b\bigr),\\
 &\forall a,b,c<n\ \bigl((E(a,b)\land E(b,c))\Longrightarrow E(a,c)\bigr),
\end{align*}
and the following two assertions for every $a,b<n$:
\begin{align*}
 \exists u<n\ \bigl(&E(a,u)\land E(b,u)\land
   \forall v<n\,((E(a,v)\land E(b,v))\Longrightarrow E(u,v))\bigr),\\
 \exists l<n\ \bigl(&E(l,a)\land E(l,b)\land
   \forall v<n\,((E(v,a)\land E(v,b))\Longrightarrow E(v,l))\bigr).
\end{align*}
Repeated bounded variables denote successive quantifiers with the
displayed bound.  Implication, equivalence, and non-strict comparison
throughout this section abbreviate their expressions using the
primitive connectives and $=,<$.
\end{definition}

The next definition gives a bounded formula for the graph criterion,
including the admissible placements and the paths.  All bound variables
are renamed when abbreviations are expanded, to avoid capture.

\begin{definition}[The witness predicate]\label{q:witness-predicate}
For $j\ge1$, the formula $Q(n,t,j)$ asserts that there is
\begin{equation}\label{q:color-encoding}
 z<n^{j^2},\qquad d(p,q)=\operatorname{dig}_n(z,pj+q)\quad(p,q<j),
\end{equation}
with the following properties.

First, the colors are symmetric, bottom occurs exactly on loops, and
the triangle inequality holds.  Explicitly, set
$\operatorname{Bot}(a)\equiv\forall b<n\ E(a,b)$ and require
\begin{align}
 &\forall p,q<j\quad d(p,q)=d(q,p),\notag\\
 &\forall p,q<j\quad
        \bigl(\operatorname{Bot}(d(p,q))\longleftrightarrow p=q\bigr),
        \label{q:graph-bottom}\\
 &\forall p,q,u<j\ \forall b<n\quad
 \bigl(E(d(p,u),b)\land E(d(u,q),b)\bigr)
                       \Longrightarrow E(d(p,q),b).\notag
\end{align}
Second, require order reflection:
\begin{equation}\label{q:graph-reflect}
 \forall a,b<n\quad \neg E(a,b)\Longrightarrow
  \exists p,q<j\ \bigl(E(d(p,q),a)\land\neg E(d(p,q),b)\bigr).
\end{equation}

For the third requirement, let $0<s<2^{j^2}$ be a nonempty ordered-pair
mask and write
$\operatorname{Sel}_s(p,q)\equiv\operatorname{dig}_2(s,pj+q)=1$.
A number $f<j^j$ specifies a vertex map by
$f[p]=\operatorname{dig}_j(f,p)$.  Set
\begin{equation}\label{q:placement-arithmetic}
 \operatorname{Adm}(f)\equiv
   \forall p,q<j\ E\bigl(d(f[p],f[q]),d(p,q)\bigr).
\end{equation}
The derived undirected adjacency is the bounded predicate
\begin{equation}\label{q:adjacency-arithmetic}
 \begin{split}
 \operatorname{Adj}_s(r,r')\equiv\exists f<j^j\ \bigl(
   &\operatorname{Adm}(f)\land
       \exists a,b<j\ \bigl(\operatorname{Sel}_s(a,b)\land\\[-2pt]
   &((f[a]=r\land f[b]=r')\lor
                     (f[a]=r'\land f[b]=r))\bigr)\bigr).
 \end{split}
\end{equation}
For $h<j^{j+1}$, put $h[i]=\operatorname{dig}_j(h,i)$, and define
\begin{equation}\label{q:path-arithmetic}
 \begin{split}
 \operatorname{Path}_s(p,q)\equiv\exists h<j^{j+1}\ \bigl(
      &h[0]=p\land h[j]=q\land\\[-2pt]
      &\forall i<j\ (h[i]=h[i+1]\lor
                            \operatorname{Adj}_s(h[i],h[i+1]))\bigr).
 \end{split}
\end{equation}
Write
\begin{align}
 \operatorname{Up}_s(b)&\equiv
   \forall a,c<j\,
       \bigl(\operatorname{Sel}_s(a,c)\Longrightarrow E(d(a,c),b)\bigr),
           \label{q:mask-upper}\\
 \operatorname{BelowJoin}_s(p,q)&\equiv
   \forall b<n\,
       \bigl(\operatorname{Up}_s(b)\Longrightarrow E(d(p,q),b)\bigr).
           \notag
\end{align}
The third requirement is
\begin{equation}\label{q:graph-connect-arithmetic}
 \forall s<2^{j^2}\quad s\ne0\Longrightarrow
   \forall p,q<j\,
     \bigl(\operatorname{BelowJoin}_s(p,q)
                          \Longrightarrow\operatorname{Path}_s(p,q)\bigr).
\end{equation}
The conjunction of all three initial color requirements---symmetry,
the bottom condition \eqref{q:graph-bottom}, and the triangle condition---
with \eqref{q:graph-reflect} and \eqref{q:graph-connect-arithmetic},
under the existential quantifier in \eqref{q:color-encoding},
is $Q(n,t,j)$.
\end{definition}

Every expression just given expands to the register language's bounded
arithmetic.  All element and vertex variables have bounds $n$ and $j$;
the remaining bounds $n^{j^2}$, $2^{j^2}$, $j^j$, and $j^{j+1}$ are
terms.  In particular the definition uses neither an unbounded path
quantifier nor a quantifier over a function as a separate sort.

\begin{lemma}[Meaning of the arithmetic witness]\label{q:arithmetic-criterion}
For every table satisfying $P(n,t)$,
\begin{equation}\label{q:arithmetic-R}
 \mathcal R\bigl((\{0,\ldots,n-1\};\preceq)\bigr)
       \quad\Longleftrightarrow\quad
                     (\exists j\ge1)\ Q(n,t,j).
\end{equation}
Moreover $Q(n,t,j)$ holds exactly when the colored-graph criterion has
a witness with $j$ vertices.
\end{lemma}

\begin{proof}
For a table satisfying $P(n,t)$, the digit encoding lists all color
tables on $j$ vertices.  The predicate $\operatorname{Bot}$ identifies
the least lattice element regardless of its numerical label.  Being
below every common upper bound is equivalent to being below the join,
so the triangle clause and $\operatorname{BelowJoin}_s$ express the
corresponding lattice inequalities in \Cref{f:criterion}.

The placement codes list all selfmaps of the vertex set.  The nonzero
ordered-pair masks select exactly the same edge sets as undirected
masks: colors are symmetric, and $\operatorname{Adj}_s$ permits either
ordering of the endpoints.  Every connected pair has a path of at most
$j-1$ edges; equality steps pad it to the $j$ steps encoded by $h$.
Conversely, every encoded path certifies connectivity.  The base-one
digit convention covers $n=1$ and $j=1$, and the masks include loops.
Thus $Q(n,t,j)$ holds exactly when there is a graph witness with $j$
vertices.  The representability equivalence now follows from
\Cref{f:criterion}.
\end{proof}

\subsection{The finite table selection}

\begin{definition}[Carrier and order]\label{q:minimum-table}
Let
\[
 \mathcal I=\{(n,t):1\le n\le M_*,\ 0\le t<2^{n^2}\}
\]
have increasing pair-lexicographic order, with $n$ compared first.
Define
\begin{equation}\label{q:selection-z}
 z_{n,t}=\begin{cases}
  1,& P(n,t)\ \text{and}\ (\forall j=1,\ldots,B)\ \neg Q(n,t,j),\\
  0,&\text{otherwise},
 \end{cases}
\end{equation}
and, with the empty product equal to one,
\begin{equation}\label{q:selection-w}
        w_{n,t}=z_{n,t}
                \prod_{\substack{(m,u)\in\mathcal I\\(m,u)<(n,t)}}
                              (1-z_{m,u}).
\end{equation}
Set
\begin{equation}\label{q:a0t0}
       a_0=\sum_{(n,t)\in\mathcal I}n w_{n,t},\qquad
       t_0=\sum_{(n,t)\in\mathcal I}t w_{n,t}.
\end{equation}
The structure $L_0$ has carrier $\{i:0\le i<a_0\}$ and order
\begin{equation}\label{q:minimum-order}
               i\le_{L_0}l\quad\Longleftrightarrow\quad
                    \operatorname{dig}_2(t_0,ia_0+l)=1.
\end{equation}
\end{definition}

All products, sums, and tests in this definition are bounded finite
constructions once $B$ is fixed.  We next prove that the cutoff is
exhaustive on this entire initial segment, so that a failed bounded
test selects a negative instance.

\begin{lemma}[Size of a graph-witness search]\label{q:search-size}
For $1\le n\le M_*$ and $0\le t<2^{n^2}$, there is a program of
\Cref{q:register-language}, of size less than $k$, which terminates
exactly when $Q(n,t,j)$ holds for some positive $j$.  On termination
it outputs the least such $j$.  The program contains neither $B$ nor
any oracle instruction.
\end{lemma}

\begin{proof}
Fix this particular pair $n,t$.  Consider the following program schema,
where the first two right-hand sides are arithmetic terms for the
specified binary numerals and $Q$ is expanded as in
\Cref{q:witness-predicate}:
\[
 \begin{array}{l}
 v_1:=\text{the numeral for }n;\\
 v_2:=\text{the numeral for }t;\\
 v_3:=1;\\
 \textbf{while }\neg Q(v_1,v_2,v_3)\textbf{ do }v_3:=v_3+1;\\
 v_0:=v_3.
 \end{array}
\]
Binary sequencing makes this a program of
\Cref{q:register-language}.  Each test of $Q$ is bounded and therefore
terminates.  The loop tests the positive integers in order, proving the
claimed termination and output properties without assuming $P(n,t)$.
Writing the numerals by binary Horner
expressions uses only $0,1$, addition, and multiplication.  Since
$t<2^{n^2}$, the two numeral terms cost at most $100(n+1)^3$ nodes,
with ample allowance for their fixed register targets.

The bounded formula templates and the rest of the program schema are
fixed independently of $n,t$.  Their expansion can be counted directly.
Use $2=(1+1)$, replace $A\Rightarrow B$ by $\neg A\lor B$, replace
$A\Leftrightarrow B$ by
$(\neg A\lor B)\land(\neg B\lor A)$, and use binary conjunctions
and one binder for each successive quantifier.  Substitute every
abbreviation separately, without shared subexpressions or loop
unrolling.  Counting a variable occurrence initially as one gives:
\begin{center}
\begin{tabular}{lrrr}
\toprule
Clause & Nodes & Variable occurrences & Binders\\
\midrule
Symmetry & 53 & 26 & 2\\
Bottom condition & 173 & 68 & 4\\
Triangle condition & 240 & 92 & 4\\
Order reflection & 198 & 72 & 4\\
Mask and connectivity & 717 & 281 & 13\\
\midrule
Full $Q$ & 1394 & 542 & 28\\
\bottomrule
\end{tabular}
\end{center}
The last line adds four conjunction nodes and the outer existential:
its binder and the seven-node bound $n^{j^{1+1}}$ contribute nine
nodes with the quantifier itself.  Renaming the expanded binders
consecutively as $v_{100},\ldots,v_{127}$ avoids all capture; the
bound of each quantifier is renamed in the preceding environment.
Only $v_1,v_2,v_3$ remain free.  Every node then costs at most $128$,
so $Q$ costs at most $128\cdot1394=178432$.
Outside the two numeral assignments, the rest of the program costs
$29$: initialization costs $6$, the while and negation nodes cost $2$,
increment costs $11$, final output costs $6$, and the four sequencing
nodes cost $4$.  In particular the following much larger allowance
bounds the entire program:
\begin{equation}\label{q:program-size}
 100(n+1)^3+(2\cdot10^6)^{204}
                      <2^{(M_*+10)^4}=k.
\end{equation}
The second summand exceeds $178461$.  For the strict inequality,
$M_*+1\le2^{M_*}$ bounds the first summand by $2^{3M_*+7}$,
and $2\cdot10^6<2^{21}$ bounds the second by $2^{4284}$.
Since $M_*>10^9$, their sum is less than $2^{3M_*+8}$, which
is less than $2^{(M_*+10)^4}$.

\end{proof}

\begin{lemma}[Uniform Busy Beaver witness cutoff]\label{q:witness-cutoff}
If $1\le n\le M_*$, $0\le t<2^{n^2}$, and $P(n,t)$ holds, then
\begin{equation}\label{q:finite-cutoff}
 \mathcal R(T_{n,t})\quad\Longleftrightarrow\quad
                   (\exists j\in\{1,\ldots,B\})\ Q(n,t,j).
\end{equation}
In fact every representable such table has a witness of size strictly
less than $B$.
\end{lemma}

\begin{proof}
For a representable table, \Cref{q:arithmetic-criterion} makes the
program in \Cref{q:search-size} terminate with a witness size $j$.
Its size is less than $k$, so the definition of $B$ gives $j<B$.
Conversely, any $j\le B$ satisfying $Q(n,t,j)$ supplies a
representation by \Cref{q:arithmetic-criterion}.
\end{proof}

\begin{theorem}[Finite description of a minimum example]\label{q:minimum}
The bounded formulas \eqref{q:selection-z}--\eqref{q:minimum-order},
with the ordinary unevaluated Busy Beaver constant
\eqref{q:BB-bound}, specify the full finite order table of a
nonrepresentable lattice $L_0$ of minimum cardinality.  Its cardinality
is the integer $a_0$ in \eqref{q:a0t0}, and
\[
                              7<a_0\le M_*.
\]
This is a finite order description and a proof of minimum cardinality;
it does not supply a numerically evaluated value of that cardinality.
\end{theorem}

\begin{proof}
By \Cref{q:witness-cutoff}, for every index in $\mathcal I$ the bit
$z_{n,t}$ is one exactly when the table is a nonrepresentable lattice.
By \Cref{q:size-ceiling}, some such bit is one.  Let $(n_*,t_*)$ be
the first index for which it is one.  The product in
\eqref{q:selection-w} is one at that index; at each later index it
contains the zero factor $1-z_{n_*,t_*}$, and at each earlier index
the factor $z_{n,t}$ is zero.  Hence exactly one $w_{n,t}$ equals
one.  Formula \eqref{q:a0t0} gives $a_0=n_*$ and $t_0=t_*$, so
\eqref{q:minimum-order} is precisely this lattice's full order table.
It is nonrepresentable.

Every nonempty lattice of cardinality smaller than $a_0$ has an
order table among earlier indices in $\mathcal I$.  Its selection
bit is zero and its lattice predicate is true, so the cutoff lemma
makes it representable.  Thus $L_0$ has minimum cardinality among all
finite nonrepresentable lattices.  The represented range through
seven elements in \Cref{f:small-lattices} gives $a_0>7$, and its
index gives $a_0\le M_*$.  Nothing in this argument evaluates $B$:
all its uses concern the outputs of the terminating programs in its
definition.  The result is therefore the claimed finite description
with an exhaustive bound, rather than a numerical enumeration.
\end{proof}

The constants and entries above are unambiguous in the standard natural
numbers.  Their definition supplies no ordinary computable procedure
for evaluating the Busy Beaver constant, no evaluated decimal value
of $a_0$, and no conventional small diagram of the minimum lattice.
Those would be stronger forms of identification than the finite
description proved here.
